% !TEX root = ../../ctfp-print.tex

\lettrine[lhang=0.17]{范}{畴论书籍很难找到一个完美的结束点}。总有更多知识等待探索。范畴论是一个庞大的学科，同时我们不断发现相同主题、概念与模式反复出现。有种说法是所有概念都是Kan扩张，事实上你可以用Kan扩张推导出极限、余极限、伴随函子、单子、Yoneda引理等诸多内容。范畴概念本身出现在所有抽象层级中，幺半群和单子的概念也是如此。哪个是最基本的？事实证明它们相互关联，形成永无止境的抽象循环。我认为展示这些相互联系可能是本书的理想结尾方式。

\section{双范畴}

范畴论最困难的方面之一在于视角的持续切换。以集合范畴为例：我们习惯通过元素定义集合——空集没有元素，单元素集包含一个元素，两个集合的笛卡尔积是有序对的集合等等。但在讨论范畴$\Set$时，我要求你忘记集合的具体内容，转而关注对象间的态射（箭头）。你偶尔可以掀开面纱观察特定泛构造在$\Set$中对应的具体元素描述：终对象被揭示为单元素集等等。但这些只是合理性验证。

函子定义为范畴间的映射。将映射视为范畴中的态射是很自然的想法。函子实际上是范畴范畴（若想避免大小问题则考虑小范畴）中的态射。当我们把函子当作箭头时，就失去了它对范畴内部结构（对象和态射）作用的信息，就像在$\Set$中将函数视为箭头时失去其对元素作用的细节一样。但任意两个范畴间的函子也构成范畴。这次你需要将原本作为箭头的对象视为另一个范畴中的对象——在函子范畴中，函子是对象而自然变换是态射。我们发现同一事物可以既是某个范畴的箭头又是另一个范畴的对象，朴素的"对象为名词、箭头为动词"观点不再成立。

与其在两种视角间切换，不如尝试将其融合为一。这就得到了$\cat{2}$-范畴概念：其中对象称为$0$-胞腔，态射称为$1$-胞腔，态射间的态射称为$2$-胞腔。

\begin{figure}[H]
  \centering
  \includegraphics[width=0.35\textwidth]{images/twocat.png}
  \caption{$0$-胞腔 $a, b$；$1$-胞腔 $f, g$；以及$2$-胞腔 $\alpha$。}
\end{figure}

\noindent
范畴的范畴$\Cat$就是典型示例：$0$-胞腔是范畴，$1$-胞腔是函子，$2$-胞腔是自然变换。$\cat{2}$-范畴定律指出任意两个$0$-胞腔间的$1$-胞腔构成范畴（即$\cat{C}(a, b)$是同态范畴而非同态集）。这完美契合我们早先关于任意两个范畴间函子构成函子范畴的论述。

特别地，从任一$0$-胞腔回到自身的$1$-胞腔也构成范畴$\cat{C}(a, a)$，但该范畴具有更强结构：$\cat{C}(a, a)$成员既可视为$\cat{C}$中的箭头，也可视为$\cat{C}(a, a)$中的对象。作为箭头它们可以相互组合，但当视为对象时，组合就变成了从对象对到对象的映射。实际上这非常类似于张量积——确切地说，这种张量积以恒等$1$-胞腔为单位元。结果表明，在任意$\cat{2}$-范畴中，同态范畴$\cat{C}(a, a)$自动成为单子范畴，其张量积定义为$1$-胞腔的组合。结合律和单位律直接继承自范畴相应的公理。

让我们回到典范示例$\Cat$来具体分析：同态范畴$\Cat(a, a)$即$a$上的自函子范畴。自函子组合在此扮演张量积角色，恒等函子是该乘积的单位元。我们之前见过自函子构成单子范畴（在定义单子时利用过此性质），但现在发现这是更普遍现象：任意$\cat{2}$-范畴中的自$1$-胞腔都构成单子范畴。稍后推广单子概念时会再次涉及这点。

或许记得在一般单子范畴中，我们不要求幺半群定律严格成立——通常只需满足到同构即可。在$\cat{2}$-范畴中，$\cat{C}(a, a)$的单子定律由$1$-胞腔组合定律直接导出，这些定律是严格的，因此总会得到严格单子范畴。不过也可以进一步放松这些定律：例如恒等$1$-胞腔$\idarrow[a]$与另一$1$-胞腔$f \Colon a \to b$的组合可以仅仅是同构而非相等。$1$-胞腔的同构通过$2$-胞腔定义——换句话说，存在具有逆元的$2$-胞腔：
\[\rho \Colon f \circ \idarrow[a] \to f\]

\begin{figure}[H]
  \centering
  \includegraphics[width=0.35\textwidth]{images/bicat.png}
  \caption{双范畴中的单位律成立到同构（可逆$2$-胞腔$\rho$）。}
\end{figure}

\noindent
我们可以对左单位律和结合律做同样处理。这种放松版$\cat{2}$-范畴称为双范畴（此处省略了一些附加协调性定律）。

正如预期，双范畴中的自$1$-胞腔构成具非严格定律的一般单子范畴。

双范畴的一个有趣例子是span范畴。对象$a$与$b$间的span由对象$x$及一对态射组成：
\begin{gather*}
  f \Colon x \to a \\
  g \Colon x \to b
\end{gather*}

\begin{figure}[H]
  \centering
  \includegraphics[width=0.35\textwidth]{images/span.png}
\end{figure}

\noindent
你可能记得我们在定义范畴积时使用过span。这里我们要将span视为双范畴中的$1$-胞腔。第一步是定义span的组合：假设存在相邻span：
\begin{gather*}
  f' \Colon y \to b \\
  g' \Colon y \to c
\end{gather*}

\begin{figure}[H]
  \centering
  \includegraphics[width=0.5\textwidth]{images/compspan.png}
\end{figure}

\noindent
组合将产生第三个span，其顶点为$z$。最自然的选择是$g$沿$f'$的拉回。记住拉回是由对象$z$及两个态射组成的：
\begin{align*}
  h  & \Colon z \to x \\
  h' & \Colon z \to y
\end{align*}
使得：
\[g \circ h = f' \circ h'\]
并且对所有此类对象具有泛性质。

\begin{figure}[H]
  \centering
  \includegraphics[width=0.5\textwidth]{images/pullspan.png}
\end{figure}

\noindent
暂时专注于集合范畴上的span：此时拉回就是笛卡尔积$x \times y$中满足
\[g\ p = f'\ q\]
的所有有序对$(p, q)$构成的集合。共享相同端点的两个span间的态射定义为顶点间的态射$h$，使得相应三角图交换。

\begin{figure}[H]
  \centering
  \includegraphics[width=0.4\textwidth]{images/morphspan.png}
  \caption{$\cat{Span}$中的$2$-胞腔。}
\end{figure}

\noindent
总结：在双范畴$\cat{Span}$中，$0$-胞腔是集合，$1$-胞腔是span，$2$-胞腔是span态射。恒等$1$-胞腔是退化span——三个对象相同且两个态射均为恒等的情形。

此前我们见过另一个双范畴示例：\hyperref[ends-and-coends]{profunctors}的双范畴$\cat{Prof}$，其中$0$-胞腔是范畴，$1$-胞腔是协变函子，$2$-胞腔是自然变换。协变函子的组合通过余积分给出。

\section{单子}

此时你应该对单子的定义相当熟悉了：它是自函子范畴中的一个幺半群。让我们用新的视角重新审视这个定义——即自函子范畴不过是双范畴 $\Cat$ 中端 1-胞的同范畴之一小部分。我们知道它是一个单子范畴：张量积由自函子的复合给出。一个幺半群在单子范畴中被定义为一个对象——这里它是一个自函子 $T$ ——连同两个态射。自函子之间的态射就是自然变换。其中一个态射将单子单位元——恒等自函子——映射到 $T$：
$$\eta \Colon I \to T$$
第二个态射将 $T \otimes T$ 的张量积映射到 $T$。由于张量积由自函子复合给出，我们得到：
$$\mu \Colon T \circ T \to T$$

\begin{figure}[H]
  \centering
  \includegraphics[width=0.3\textwidth]{images/monad.png}
\end{figure}

\noindent
我们将这两个操作识别为定义单子的基本运算（它们在 Haskell 中被称为 \code{return} 和 \code{join}），并且知道幺半群定律会转化为单子定律。

现在从这个定义中移除所有关于自函子的提及。我们从一个双范畴 $\cat{C}$ 开始，在其中选取一个 0-胞 $a$。正如先前所见，同范畴 $\cat{C}(a, a)$ 是一个单子范畴。因此我们可以在 $\cat{C}(a, a)$ 中定义一个幺半群：选取一个 1-胞 $T$ 和两个 2-胞：
\begin{align*}
  \eta & \Colon I \to T         \\
  \mu  & \Colon T \circ T \to T
\end{align*}
使其满足幺半群定律。我们将\emph{这个}称为单子。

\begin{figure}[H]
  \centering
  \includegraphics[width=0.3\textwidth]{images/bimonad.png}
\end{figure}

\noindent
这是通过 0-胞、1-胞和 2-胞给出的更一般化的单子定义。当应用于双范畴 $\Cat$ 时，它就退化为通常的单子定义。但现在让我们看看在其他双范畴中的情形。

让我们在 $\cat{Span}$ 中构造一个单子。我们选取一个 0-胞——出于即将清晰的原因，将其命名为 $\mathit{Ob}$。接着选取一个端 1-胞：一个从 $\mathit{Ob}$ 到自身的纺锤形。其顶端有一个集合，记作 $\mathit{Ar}$，配备两个函数：
\begin{align*}
  \mathit{dom} & \Colon \mathit{Ar} \to \mathit{Ob} \\
  \mathit{cod} & \Colon \mathit{Ar} \to \mathit{Ob}
\end{align*}

\begin{figure}[H]
  \centering
  \includegraphics[width=0.3\textwidth]{images/spanmonad.png}
\end{figure}

\noindent
让我们将集合 $\mathit{Ar}$ 的元素称为「箭头」。如果我还告诉你将 $\mathit{Ob}$ 的元素称为「对象」，你可能会猜到这指向何方。两个函数 $\mathit{dom}$ 和 $\mathit{cod}$ 分别为每个「箭头」指派定义域和上定义域。

要使我们的纺锤成为单子，需要两个 2-胞 $\eta$ 和 $\mu$。在此情形下单子单位元是一个平凡纺锤：从 $\mathit{Ob}$ 到 $\mathit{Ob}$ 顶端在 $\mathit{Ob}$ 的纺锤，配有两个恒等函数。2-胞 $\eta$ 是顶端 $\mathit{Ob}$ 和 $\mathit{Ar}$ 之间的函数。换句话说，$\eta$ 为每个「对象」指派一个「箭头」。$\cat{Span}$ 中的 2-胞必须满足交换条件——在此情形下：
\begin{align*}
  \mathit{dom} & \circ \eta = \id \\
  \mathit{cod} & \circ \eta = \id
\end{align*}

\begin{figure}[H]
  \centering
  \includegraphics[width=0.4\textwidth]{images/spanunit.png}
\end{figure}

\noindent
用分量表示就是：
$$\mathit{dom}\ (\eta\ \mathit{ob}) = \mathit{ob} = \mathit{cod}\ (\eta\ \mathit{ob})$$
其中 $\mathit{ob}$ 是 $\mathit{Ob}$ 中的一个「对象」。换句话说，$\eta$ 为每个「对象」指派一个以该对象为定义域和上定义域的「箭头」。我们将这种特殊的「箭头」称为「恒等箭头」。

第二个 2-胞 $\mu$ 作用于纺锤 $\mathit{Ar}$ 与自身的复合。该复合定义为拉回，因此其元素是来自 $\mathit{Ar}$ 的元素对——即「箭头」对 $(a_1, a_2)$。拉回条件为：
$$\mathit{cod}\ a_1 = \mathit{dom}\ a_2$$
我们称 $a_1$ 和 $a_2$ 是「可复合的」，因为其中一个的定义域恰是另一个的上定义域。

\begin{figure}[H]
  \centering
  \includegraphics[width=0.5\textwidth]{images/spanmul.png}
\end{figure}

\noindent
2-胞 $\mu$ 是一个函数，将一对可复合箭头 $(a_1, a_2)$ 映射为 $\mathit{Ar}$ 中的单个箭头 $a_3$。换句话说 $\mu$ 定义了箭头的复合。

容易验证单子定律对应于箭头的恒等律和结合律。我们刚刚定义了一个范畴（确切地说是一个小范畴，其中对象和箭头构成集合）。

综上所述，一个范畴本质上就是纺锤双范畴中的一个单子。

这个结果令人惊叹之处在于它将范畴与其他代数结构（如单子和幺半群）置于同等地位。范畴本身并无特殊之处——它只是两个集合和四个函数而已。事实上我们甚至不需要单独的对象集合，因为对象可以通过恒等箭头来标识（它们之间存在一一对应）。所以本质上它只是一个集合和若干函数而已。考虑到范畴论在整个数学中的核心地位，这是一个令人谦卑的认识。

\section{挑战}

\begin{enumerate}
  \tightlist
  \item
        推导定义为双范畴中自同伦(endomorphic) $1$-胞腔复合的张量积的单位元律与结合律。
  \item
        验证范畴 $\cat{Span}$ 中单子(monad)的公理对应所得范畴中的恒等律与结合律。
  \item
        证明范畴 $\cat{Prof}$ 中的单子是对象恒等(identity-on-objects)函子。
  \item
        在范畴 $\cat{Span}$ 中，单子的代数(mondad algebra)是什么？
\end{enumerate}

\section{参考文献}
\begin{enumerate}
  \tightlist
  \item
        \urlref{https://graphicallinearalgebra.net/2017/04/16/a-monoid-is-a-category-a-category-is-a-monad-a-monad-is-a-monoid/}{Paweł Sobociński的博客（Paweł Sobociński’s blog）}.
\end{enumerate}